% =============================================================================
%  main.tex - preambulo, portada y \input de las secciones.
%
%  ESTE ARCHIVO ES DE TODO EL EQUIPO. Casi nunca debe cambiar: solo se toca
%  para agregar una seccion nueva al final (un \input mas).
%  Tu texto va en secciones/<tu-seccion>.tex.
%
%  Compilar desde esta carpeta:   latexmk        (el .latexmkrc ya hace el resto)
% =============================================================================

% nonatbib: elsarticle carga natbib por default y natbib es INCOMPATIBLE con
% biblatex. Esta opcion le dice a la clase que no lo cargue.
% 5p = formato Elsevier a dos columnas. Si lo quieren a una columna, cambien
% 5p por 3p y no hay que tocar nada mas.
\documentclass[5p,nonatbib]{elsarticle}

% --- Quitar el pie de pagina "Preprint submitted to Elsevier" -----------------
% No estamos mandando esto a una revista, es un reporte de clase.
\makeatletter
\def\ps@pprintTitle{%
  \let\@oddhead\@empty
  \let\@evenhead\@empty
  \let\@oddfoot\@empty
  \let\@evenfoot\@empty}
\makeatother

% --- Codificacion y tipografia -----------------------------------------------
\usepackage[T1]{fontenc}      % acentos y "ñ" como un solo glifo (permite copiar
                              % texto del PDF y que la division silabica sirva)
\usepackage[utf8]{inputenc}   % el .tex se escribe en UTF-8
                              % Sin paquete de tipografia: se usa la default de
                              % LaTeX (Computer Modern). Ojo: NO usar la opcion
                              % [times] de la clase, que la cambiaria.

\usepackage{microtype}        % micro-ajustes de espaciado; menos renglones feos

% A dos columnas las columnas son angostas y una palabra larga sin guiones
% (una ruta como references.bib, o un \verb) se sale del margen. Esto deja que
% LaTeX estire un poco los espacios antes de rendirse y desbordar.
\emergencystretch=2em

% --- Idioma ------------------------------------------------------------------
% Los ajustes de babel-spanish NO se pueden poner como opciones de \usepackage:
% babel leeria "es-tabla" como si fuera un idioma y truena. Van aqui, en
% \spanishoptions, ANTES de cargar babel.
%   es-tabla        -> "Tabla 1" en vez de "Cuadro 1"
%   es-nodecimaldot -> deja el punto decimal en matematicas (12.46 y no 12,46),
%                      porque asi salen los numeros de pandas/R
%   es-noquoting    -> evita que babel se pelee con otros paquetes por comillas
\def\spanishoptions{es-tabla,es-nodecimaldot,es-noquoting}
\usepackage[spanish]{babel}
\usepackage{csquotes}         % \enquote{...} -> «...» segun el idioma activo

% --- Matematicas -------------------------------------------------------------
\usepackage{amsmath}
\usepackage{amssymb}
\usepackage{amsthm}

% Entornos tipo teorema. El estilo "definition" deja el cuerpo en letra recta
% (el default "plain" lo pone en italicas, que es para teoremas).
% Uso: \begin{definition}[Nombre opcional] ... \end{definition}
\theoremstyle{definition}
\newtheorem{definition}{Definición}

% --- Tablas y figuras --------------------------------------------------------
\usepackage{booktabs}         % \toprule \midrule \bottomrule (tablas sin lineas
                              % verticales, que es el estilo correcto)
\usepackage{tabularx}         % columnas X que se reparten el ancho sobrante
\usepackage{ragged2e}         % \RaggedRight: texto sin justificar en celdas
\usepackage{multirow}
\usepackage{graphicx}
\usepackage{multicol}
\graphicspath{{figuras/}}     % con esto basta \includegraphics{archivo.pdf}

% A dos columnas, un table* [t] cerca del final del documento puede empujar
% el resto del texto a una pagina nueva; si lo que queda despues no llena
% esa pagina, una columna se queda vacia. flushend balancea las dos columnas
% de la ULTIMA pagina para que ese sobrante no se vea como espacio en blanco.
\usepackage{flushend}
\usepackage{blindtext}


% El reporte tiene muchas tablas grandes (sobre todo en el diccionario de
% datos); el espaciado por defecto de LaTeX alrededor de flotantes y listas
% deja demasiado blanco. Se aprieta aqui para que el documento se vea mas
% denso sin tocar el tamano de fuente ni los margenes de la clase.
\setlength{\intextsep}{6pt plus 2pt minus 2pt}
\setlength{\textfloatsep}{8pt plus 2pt minus 2pt}
\setlength{\floatsep}{6pt plus 2pt minus 2pt}
\setlength{\abovecaptionskip}{4pt}
\setlength{\belowcaptionskip}{2pt}

% Parametros de colocacion de flotantes. Los defaults de LaTeX son
% conservadores y en este reporte se quedan cortos: la etapa 2 aporta 24
% flotantes contra unas 5 paginas de prosa, y como cada uno solo puede
% aparecer en o despues del punto donde se cita, la cola se desbordaba al
% final del documento (seis figuras quedaban DESPUES de la bibliografia).
%
% Lo que estorbaba, en concreto:
%   \textfraction 0.2 -> exigia que al menos 20%% de cada pagina fuera texto
%   \topnumber    2   -> maximo dos flotantes en la zona superior
%   \floatpagefraction 0.5 -> una pagina solo de flotantes tenia que ir medio
%                             llena, si no LaTeX prefiere seguir difiriendolos
%
% Esto NO recorta contenido ni cambia el tamano de letra: solo autoriza a
% LaTeX a empaquetar los flotantes mas densamente. Bajo el documento de 13 a
% 12 paginas y dejo la bibliografia al final, que es donde debe estar.
% Si en algun momento el reporte se ve demasiado saturado de figuras por
% pagina, estos son los diales que hay que aflojar (subir \textfraction,
% bajar \topnumber).
\renewcommand{\topfraction}{0.92}
\renewcommand{\bottomfraction}{0.6}
\renewcommand{\textfraction}{0.06}
\renewcommand{\floatpagefraction}{0.55}
\renewcommand{\dbltopfraction}{0.92}      %% idem para flotantes de ancho completo
\renewcommand{\dblfloatpagefraction}{0.55}
\setcounter{topnumber}{4}
\setcounter{bottomnumber}{3}
\setcounter{totalnumber}{6}
\setcounter{dbltopnumber}{3}

\usepackage{enumitem}
\setlist{noitemsep, topsep=3pt, partopsep=0pt}

% --- Bibliografia (biblatex + biber) -----------------------------------------
% style=authoryear -> (Jolliffe, 2002) en el texto.
% Para citas numericas [1], cambien authoryear por numeric-comp.
\usepackage[backend=biber,style=apa,sorting=nyt,maxbibnames=99]{biblatex}
% La bibliografia se compone un punto mas chica que el cuerpo, convencion
% habitual en revistas: son datos de consulta, no prosa que se lea seguida.
% Con las URLs largas de las fuentes normativas esto ahorra ~10 renglones,
% que es lo que hacia que se derramara a una pagina extra.
\renewcommand*{\bibfont}{\small}
\setlength{\bibitemsep}{0pt}      %% sin aire extra entre entradas
\setlength{\bibhang}{1em}

% --- Referencias cruzadas (SIEMPRE hasta el final) ---------------------------
\usepackage[colorlinks=true, linkcolor=blue,citecolor=blue]{hyperref}
\usepackage[spanish]{cleveref} % \cref{fig:x} escribe solo "Figura 1"

\usepackage[version=4]{mhchem} % formulas de quimica

% cleveref en español dice "apartado 2"; en un reporte queremos "Sección 2".
% \cref -> minuscula (a media oracion), \Cref -> mayuscula (al inicio).
\crefname{section}{sección}{secciones}
\Crefname{section}{Sección}{Secciones}
\crefname{table}{tabla}{tablas}
\Crefname{table}{Tabla}{Tablas}
\crefname{equation}{ecuación}{ecuaciones}
\Crefname{equation}{Ecuación}{Ecuaciones}

% --- Textos de la portada en español ----------------------------------------
% elsarticle trae "Abstract" y "Keywords" escritos a fuerza en ingles;
% la clase deja estos dos ganchos para cambiarlos.
\abstracttitle{Resumen}
\keywordtitle{Palabras clave}


% BIBLIOGRAFIAS AQUI
\addbibresource{secciones/introduccion.bib}
\addbibresource{secciones/revision.bib}
\addbibresource{secciones/incertidumbre.bib}

\begin{document}

% --- Carátula ----------------------------------------------------------------
% La rúbrica pide número de equipo, nombre y matrícula de cada integrante y el
% nombre del proyecto vinculado con el objetivo. elsarticle no trae carátula,
% así que se compone a mano antes del frontmatter (que pone el resumen)

\begin{frontmatter}

\title{Etapa 1 de la iniciativa de investigación del uso de modelación estocástica para la reforestación}

\author[tec]{Angel Eduardo-Luna}
\author[tec]{Irasema Alvarez}
\author[tec]{Carolina Sánchez Meneses}
\author[tec]{Andrea Estefanía de la Peña}
\address[tec]{Instituto Tecnológico y de Estudios Superiores de Monterrey, MA2004B--- Grupo 202, equipo 5}

% Resumen breve de la etapa 1 (recoleccion de informacion): problema,
% informacion reunida y siguiente paso. Sin citas, formulas ni referencias.
\begin{abstract}
La Comisión Nacional Forestal planea reforestar 30 polígonos en el matorral
desértico de San Luis Potosí con diez especies nativas, plantadas al
tresbolillo con una meta de 658 plantas por hectárea. El reto plantea dos
decisiones: distribuir las especies dentro de cada polígono para reducir la
competencia entre plantas vecinas y estimar cuántas plantas de cada especie
comprar cuando el número de plantas existentes y la supervivencia de las
nuevas son inciertos.

Este documento corresponde a la primera etapa del proyecto, dedicada a
reunir la información necesaria para plantear el modelo. Se describe el
contexto de la reforestación y de la competencia entre plantas en zonas
áridas, se identifican las variables de decisión, los parámetros, las
fuentes de incertidumbre y la información que aún falta obtener, y se
revisan trabajos que tratan la asignación espacial como un problema de
coloración de grafos con penalización por conflicto. También se reúnen los
conceptos de probabilidad, simulación Monte Carlo y cadenas de Markov con
los que se estimará la cantidad de plantas a adquirir, y se esboza una
primera identificación del objetivo y de las restricciones del modelo de
optimización. La formulación completa, su solución y su validación
corresponden a las etapas siguientes.
\end{abstract}

\begin{keyword}
reforestación \sep tresbolillo \sep competencia interespecífica \sep
coloración de grafos \sep simulación Monte Carlo \sep cadenas de Markov
\end{keyword}
\end{frontmatter}

% --- Secciones ---------------------------------------------------------------
% Una seccion = un archivo = una issue = una branch.
% El orden sigue la rubrica del curso: introduccion y justificacion,
% problematica y objetivo, preparacion de la base de datos, modelacion y
% validacion, resultados, discusion y conclusiones. Cada archivo de
% secciones/ abre con el nivel de encabezado que le toca aqui; al mover un
% archivo de lugar hay que revisar ese nivel.

\section{Introducción y contexto del problema}
\label{sec:introduccion}
 
\subsection{La reforestación y su objetivo}
La deforestación por cambio de uso de suelo, impulsada sobre todo por la expansión agrícola, la ganadería y la infraestructura, es una de las principales causas de pérdida de biodiversidad y de servicios ecosistémicos \autocite{geist2002}. Una de las respuestas a este problema es la restauración ecológica, definida como el proceso de asistir la recuperación de un ecosistema que ha sido degradado, dañado o destruido \autocite{ser2004}.
La reforestación es una de las herramientas de la restauración, la cual consiste en establecer vegetación, mediante siembra o plantación, en terrenos que antes la tenían y la perdieron \autocite{conafor2010}. Su objetivo no es solo recuperar la cobertura vegetal, sino también los servicios que el ecosistema presta, como la retención de suelo, la infiltración de agua, el hábitat para la fauna y la captura de carbono \autocite{chazdon2008}.
 
\subsection{Tipos y estrategias de reforestación}
Según el propósito, la reforestación puede ser de conservación o restauración, comercial, agroforestal o urbana \autocite{conafor2010}. Según el grado de intervención, se distingue entre la restauración pasiva, que retira el disturbio y deja que el sitio se regenere de forma natural, y la activa, la cual planta especies cuando la regeneración natural no es suficiente \autocite{holl2011}. Entre ambas existen estrategias intermedias como la nucleación, que consiste en crear pequeños núcleos (islas de árboles, perchas para aves, transposición de suelo) que aceleran la llegada de semillas y la regeneración alrededor de ellos \autocite{reis2010,corbin2012}. La elección depende del grado y tipo de disturbio y de los objetivos del proyecto \autocite{ser2004}.
 
\subsection{Cómo se realiza una reforestación}
En la mayoría de los casos, una reforestación sigue tres etapas \autocite{conafor2010}. La planeación incluye el diagnóstico del sitio (clima, suelo, pendiente, vegetación existente y disturbios), la selección de especies, la definición de la densidad y del diseño de plantación, y la obtención de la planta en vivero. La ejecución comprende la preparación del terreno, el transporte de la planta y la plantación en la temporada de lluvias. Finalmente, el mantenimiento y seguimiento abarcan la protección contra el ganado y el fuego, la reposición de plantas muertas y la medición de la supervivencia.
 
\subsection{Selección de especies, densidad y distribución}
Las especies se eligen a partir del tipo de vegetación del área que se tiene, dando prioridad a especies nativas adaptadas al clima y al suelo local \autocite{conafor2010}. Una forma práctica de hacerlo es un levantamiento de la vegetación existente en la región para plantar las especies en proporción a su frecuencia de aparición.
La densidad, es decir, el número de plantas por hectárea, y su distribución dependen tanto del tamaño que alcanza cada especie en edad adulta, como de la disponibilidad de agua y nutrientes, de la pendiente y del objetivo de la plantación \autocite{conafor2010}. Los diseños más usados son el marco real (cuadrícula) y el tresbolillo, en el que las filas se desfasan para formar triángulos equiláteros. El tresbolillo aprovecha mejor el espacio y protege mejor contra el viento y la erosión.
 
\subsection{Competencia entre plantas}
La competencia ocurre cuando plantas vecinas usan los mismos recursos limitados (agua, luz y nutrientes del suelo) y el consumo de una reduce lo disponible para la otra planta \autocite{craine2013}. Es más intensa entre plantas que están cerca, que son de tamaño parecido y que toman los recursos de la misma forma, a la misma profundidad de raíz y en la misma época.
Sin embargo, en zonas áridas las interacciones no son solo de competencia. Algunas plantas, conocidas como nodrizas, crean bajo su copa un microambiente con más sombra, humedad y nutrientes que favorece el establecimiento de otras especies \autocite{flores2003}. El mezquite (\textit{Prosopis}) es una nodriza común en el norte de México, y aprovechar este efecto es una estrategia reconocida para restaurar ambientes degradados \autocite{padilla2006}.
 
\subsection{Especies invasoras y monocultivos}
Una especie invasora es una especie exótica que se establece, se reproduce y se dispersa fuera de su área natural. Puede desplazar a las especies nativas, alterar el régimen de fuego, el ciclo de nutrientes y el uso del agua, y su control suele ser muy costoso \autocite{mack2000}. Introducirla en una reforestación, por ejemplo al usar especies de rápido crecimiento ajenas al sitio, puede dañar más el ecosistema de lo que se pretendía restaurar.
Un monocultivo es una plantación de una sola especie. Aunque es más sencillo de manejar, reduce la biodiversidad, hace a la plantación más vulnerable a plagas, enfermedades y eventos climáticos extremos y presta menos servicios ecosistémicos que una plantación mixta \autocite{liu2018}.
 
\subsection{Disponibilidad de agua}
En zonas áridas y semiáridas el agua es el recurso más limitante, por lo que condiciona tanto la selección como la distribución de las especies. Conviene elegir especies tolerantes a la sequía, como agaves, nopales y yucas, y plantarlas al inicio de la temporada de lluvias. En la distribución, plantar juntas especies con alta demanda de agua aumenta la competencia, mientras que colocar especies sensibles bajo nodrizas mejora su supervivencia \autocite{padilla2006}.
 
\subsection{Problemas de una mala planeación}
Una reforestación mal planeada o mal ejecutada puede introducir especies invasoras, generar monocultivos o acomodar las plantas de forma que compitan entre sí, lo que provoca baja supervivencia y deja sin cumplir el objetivo inicial \autocite{conafor2010,holl2011}. A esto se suma un problema logístico: si se compran plantas de más se gasta dinero de forma innecesaria y muchas mueren, y si se compran de menos no se alcanza la meta de plantación.
 
\subsection{Introducción al problema del reto}
El reto se ubica en la zona de la Línea de Transmisión Eléctrica Dominica--Charcas, en San Luis Potosí, dentro de la subprovincia Llanos y Sierras Potosino--Zacatecanos. La vegetación es principalmente matorral desértico rosetófilo y micrófilo, con un deterioro notable por sobrepastoreo de ganado. El socio formador (CONAFOR) planea reforestar 30 polígonos de entre 1.28 y 8~ha con 10 especies nativas de agaves, nopales, mezquite y yuca, en un diseño de tresbolillo elegido por los fuertes vientos de la zona y con una meta de 658 plantas por hectárea.
El problema tiene dos partes. La primera es decidir cómo distribuir las especies dentro de cada polígono para reducir la competencia entre plantas vecinas, respetando la cantidad requerida por especie. La segunda es estimar cuántas plantas de cada especie se deben comprar, considerando que en cada polígono ya existen plantas que nacieron de forma natural y cuyo número no se conoce con certeza.
 
\subsection{Justificación}
Una buena distribución de especies reduce la competencia y aprovecha la facilitación entre plantas, lo que aumenta la supervivencia y con ello el éxito de la reforestación. Esto es especialmente importante en una zona árida, donde cada planta perdida es difícil de reponer. Por otro lado, estimar bien la cantidad a comprar evita tanto el gasto innecesario como el incumplimiento de la meta. Contar con un método sistemático para ambas decisiones permite a CONAFOR planear con datos en lugar de depender de un orden de plantación fijo.
 
\subsection{Objetivo general}
El objetivo es proponer un modelo matemático y computacional que permita a CONAFOR distribuir las especies dentro de cada polígono minimizando la competencia entre ellas y cumpliendo la demanda por especie, y estimar mediante simulación la cantidad de plantas por especie que conviene adquirir considerando la incertidumbre sobre las plantas ya existentes.
\section{Definición del problema}
\label{sec:definicion-problema}

\subsection{Descripción del problema}
\label{subsec:descripcion-problema}

\textbf{Qué decisiones deben tomarse.} Son dos decisiones relacionadas pero distintas.

\begin{enumerate}
\item \textbf{Distribución de especies.} En qué posición de la retícula al tresbolillo se coloca cada planta, dentro del margen que el plan permite variar por especie (de $\pm$5\% a $\pm$10\%), buscando reducir la competencia entre vecinos.
\item \textbf{Cantidad a adquirir.} Cuántas plantas de cada una de las 10 especies comprar por polígono, dado que ya existen algunas plantas cuyo número exacto no se conoce con certeza, y dada la mortalidad esperada de las plantas nuevas.
\end{enumerate}

\textbf{Qué información está disponible.} El plan de siembra (plantas por hectárea y por especie); conteos de plantas existentes por especie en cada uno de los 30 polígonos; la superficie de cada polígono; datos de supervivencia de 2021 a 2023 por especie, medidos en una muestra del 5\% de la superficie; y la secuencia de 22 posiciones que actualmente se usa para ordenar la plantación en cada 100 m lineales.

\textbf{Qué resultados se esperan obtener.} Una cantidad de plantas a comprar por especie (y, si se desglosa, por polígono) que cumpla la meta del plan considerando la incertidumbre sobre las existentes y la mortalidad esperada; y un algoritmo de fácil implementación en campo que indique en qué orden o posición plantar cada especie para reducir la competencia frente al arreglo actual.

\subsection{Variables y parámetros}
\label{subsec:variables-parametros}

\Cref{tab:variables-parametros} resume las variables de decisión, los parámetros conocidos, las variables aleatorias y la información aún no disponible que intervienen en el problema.

\begin{table*}[t]
\centering
\caption{Variables y parámetros identificados para el problema de reforestación.}
\label{tab:variables-parametros}
\small
\begin{tabular}{@{}p{0.27\textwidth}p{0.18\textwidth}p{0.47\textwidth}@{}}
\toprule
Elemento & Tipo & Descripción \\
\midrule
Especie a plantar en cada posición & Variable de decisión & Cuál de las 10 especies ocupa cada sitio de la retícula \\
Cantidad a comprar por especie & Variable de decisión & Número de plantas adicionales a adquirir \\
Plantas requeridas por especie (plan) & Parámetro conocido & 658 plantas/ha repartidas en 10 especies, según el documento del reto \\
Presupuesto & Parámetro conocido & \$13{,}200 (unidad y alcance por confirmar con el socio) \\
Distancia entre plantas & Parámetro conocido & Definida por el arreglo al tresbolillo y la densidad del plan \\
Plantas existentes por especie y polígono & Variable aleatoria & Se conoce un conteo, pero con incertidumbre sobre cuántas hay realmente en el área no muestreada \\
Supervivencia por especie & Variable aleatoria & Estimada de los datos históricos de 2021 a 2023, sujeta al tamaño de la muestra de monitoreo \\
Compatibilidad o competencia entre especies & Información a estimar & No viene en los datos del reto; debe construirse con literatura o criterio experto, y es el parámetro más incierto del modelo \\
Requerimientos de agua por especie & Información a obtener & Relevante para el matorral desértico, no incluida en los datos actuales \\
Disponibilidad de plantas en vivero & Información a obtener & Relevante para saber si la compra calculada es factible de adquirir \\
Costo de adquisición por especie & Información a obtener & Necesario para optimizar el presupuesto, no incluido en los datos actuales \\
Características del terreno (pendiente, suelo, exposición) & Información a obtener & El reto distingue tipos de vegetación y relieve por polígono, pero no se dispone del detalle por polígono \\
\bottomrule
\end{tabular}
\end{table*}

\textbf{Justificación de relevancia.} Las variables de decisión (especie por posición y cantidad por comprar) son el núcleo de las dos preguntas del reto. Las variables aleatorias (existentes y supervivencia) son las fuentes de incertidumbre que justifican usar simulación y no una fórmula determinista. La compatibilidad entre especies es indispensable porque es justamente lo que la distribución espacial busca reducir, aunque no esté en los datos. El agua, la disponibilidad, los costos y el terreno no son indispensables para una primera versión del modelo, pero condicionan qué tan realista es la solución final y son los primeros candidatos para refinar el modelo más adelante.

\section{Trabajo relacionado}
\label{sec:trabajo-relacionado}

Esta sección revisa trabajos que enfrentan problemas con una estructura similar a la del reto: asignar elementos a posiciones o momentos de modo que se reduzca la competencia, la interferencia o la incompatibilidad entre ellos. Primero se analizan cinco referencias académicas sobre reducción de competencia en reforestación, agricultura, telecomunicaciones y manejo forestal (\cref{sec:reduccion-competencia}). Después se revisan los problemas de coloración de grafos y su posible uso para representar la distribución de especies (\cref{sec:coloracion}).

%==================================================
\subsection{Reducción de competencia}
\label{sec:reduccion-competencia}

Para cada referencia se identifica el problema estudiado, las variables o elementos considerados, el criterio con el que se mide la competencia o incompatibilidad, el método de solución y los principales resultados.

\subsubsection{Competencia entre especies en un área de reforestación}

\textcite{elizalde2025} estudian un problema muy cercano al del reto.

\begin{description}
  \item[Problema.] Asignar especies de árboles a las posiciones de siembra de un polígono para minimizar la competencia interespecífica y cumplir la demanda de cada especie, siguiendo los lineamientos técnicos de CONAFOR/USPAE (625 plantas de diez especies por hectárea).
  \item[Variables.] El terreno se representa como un grafo $G=(V,A)$ cuyos vértices son las posiciones de siembra y cuyos arcos indican vecindad. La decisión es una variable binaria $x_{oi}$ (la especie $o$ se planta en la posición $i$) y una variable $R_{ij}$ con el nivel de competencia de cada arco. Los parámetros son la adyacencia $w(i,j)$, el factor de competencia $cop$ entre pares de especies y la demanda $d_o$ de cada especie.
  \item[Criterio de competencia.] Coeficientes de competencia por par de especies que solo se contabilizan cuando ambas ocupan posiciones vecinas; el objetivo es minimizar la suma de $R_{ij}$ sobre todo el terreno.
  \item[Método.] Un modelo de programación matemática exacto, resuelto con el solver CBC mediante PuLP, y un algoritmo genético con cromosoma entero (cada gen es una posición y su valor es la especie), selección por ruleta, cruce de tres puntos, mutación y elitismo.
  \item[Resultados.] Con instancias de 25 a 625 nodos y un límite de tres horas, el modelo exacto obtuvo soluciones factibles sin garantía de calidad (brecha infinita) entre 169 y 324 nodos, y no reportó resultados con 361, 400 y 625 nodos. El algoritmo genético encontró soluciones factibles en todos los tamaños, con un tiempo de 349.44 segundos en el mayor, y fue mejor o muy cercano al modelo exacto en casi todas las instancias comparables.
\end{description}

\subsubsection{Interferencia en la asignación de frecuencias}

\textcite{eisenblatter2002} presentan un estudio sobre la asignación de frecuencias en redes celulares, donde la interferencia desempeña el mismo papel que la competencia entre plantas.

\begin{description}
  \item[Problema.] Asignar canales de frecuencia a los transmisores de una red de telefonía celular evitando la interferencia o, en su defecto, minimizándola.
  \item[Variables.] Un grafo de interferencia con un vértice por transmisor y una arista cuando existe interferencia entre dos transmisores. Los canales actúan como colores. El modelo de minimización usa variables binarias $x_{vf}$ (el canal $f$ se asigna al transmisor $v$) y $z_{vw}$ (se viola una restricción en la arista $vw$), además de parámetros como los canales bloqueados y las distancias mínimas de separación.
  \item[Criterio de interferencia.] Un valor de interferencia por par de transmisores, normalizado entre 0 y 1, que distingue la interferencia cuando comparten canal y cuando usan canales adyacentes; una función de penalización evalúa cada par de frecuencias asignadas.
  \item[Método.] Los autores recorren la evolución de los modelos, desde la coloración de vértices hasta generalizaciones como la T-coloración y la coloración por listas, y formulan el problema de mínima interferencia como un programa lineal entero. Revisan heurísticas de coloración secuencial, DSATUR, \textit{threshold accepting} y algoritmos genéticos, así como cotas inferiores.
  \item[Resultados.] Todas las variantes son NP-difíciles y el problema de mínima interferencia lo es fuertemente. Las instancias reales llegan a miles de vértices. En un ejemplo, DSATUR seguido de una búsqueda local sencilla mejoró en 96\,\% el plan de un operador, mientras que las relajaciones lineales dan cotas inferiores débiles, con brechas grandes.
\end{description}

\subsubsection{Asignación espacio-temporal de cultivos}

\textcite{akplogan2011} abordan el plan de cultivos de una granja.

\begin{description}
  \item[Problema.] Asignar cultivos a unidades de terreno a lo largo de un horizonte de tiempo, considerando a la vez los aspectos espaciales y temporales y las decisiones del agricultor.
  \item[Variables.] El cultivo $x^{t}_{b,i}$ que ocupa la unidad $i$ del bloque $b$ en el año $t$. Los parámetros incluyen el tiempo mínimo de retorno de cada cultivo, el efecto del cultivo previo, la compatibilidad entre cultivo y tipo de suelo, los límites de agua de riego y los rangos de superficie por cultivo.
  \item[Criterio de incompatibilidad.] Restricciones duras (por ejemplo, cultivos no aptos para un suelo o repetidos antes del tiempo mínimo de retorno) y penalizaciones por el efecto del cultivo precedente. Además, penalizan las unidades aisladas para agrupar un mismo cultivo y las desviaciones de las superficies deseadas.
  \item[Método.] Un problema de satisfacción de restricciones ponderadas (WCSP) resuelto con búsqueda por ramificación y acotamiento en el solver Toulbar2.
  \item[Resultados.] En una granja virtual de 180~ha con hasta 120 unidades y nueve años de horizonte, los bloques independientes se resolvieron hasta el óptimo casi siempre en menos de un minuto (228 segundos en el caso más lento). Con bloques interdependientes, la instancia de 60 unidades tardó 2412.97 segundos y la de 120 no se cerró en 48 horas.
\end{description}

\subsubsection{Planeación robusta de cultivos bajo incertidumbre}

\textcite{liu2026} combinan la interacción entre cultivos con incertidumbre.

\begin{description}
  \item[Problema.] Planear la asignación de cultivos en varios años sobre un terreno heterogéneo, con incertidumbre en rendimientos, precios y costos, y aprovechando relaciones de complementariedad y competencia entre cultivos.
  \item[Variables.] Una variable binaria por unidad de terreno, periodo y cultivo; una matriz de adyacencia $W$ construida a partir de unidades vecinas en una cuadrícula (horizontal o verticalmente); una matriz de interacción $M$ entre cultivos; y parámetros de rendimiento, precio, costo, área y consumo de agua.
  \item[Criterio de competencia.] Los coeficientes de $M$ son positivos para pares complementarios (por ejemplo, leguminosa y cereal) y negativos para pares competitivos; el potencial de interacción de cada unidad se obtiene sumando, sobre sus vecinas, los términos $W_{ij}M_{c,c'}$ de los cultivos plantados.
  \item[Método.] Un marco de tres capas (heterogeneidad espacial, dinámica temporal y robustez) con optimización distribucionalmente robusta max-min, un conjunto de ambigüedad basado en la distancia de Wasserstein y escenarios generados mediante simulación Monte Carlo, implementado en PuLP.
  \item[Resultados.] En un caso del norte de China (54 unidades, 41 cultivos y siete años), el método obtuvo una ganancia en el peor caso de $2310.5$ ($10^{4}$ CNY), frente a $1820.1$ del modelo determinístico y $2180.3$ del robusto base, y elevó la proporción de leguminosas a 22\,\% (12\,\% y 15\,\% en los otros dos). Superó al robusto base en todo el rango de incertidumbre evaluado.
\end{description}

\subsubsection{Definición de competidores en bosques mixtos}

\textcite{kobal2025} estudian cómo medir la competencia entre árboles en campo.

\begin{description}
  \item[Problema.] Determinar qué árboles vecinos deben considerarse competidores para explicar el crecimiento de abetos plateados dominantes en bosques mixtos de edades diversas de las montañas dináricas (Eslovenia).
  \item[Variables.] El incremento de volumen en cinco años de 60 árboles, su volumen, y el diámetro a la altura del pecho (DBH), la altura y la distancia de 4\,454 vecinos medidos en un radio de 25.23~m. Se definen umbrales dinámicos: un vecino es competidor si su DBH supera $\alpha\cdot\mathrm{DBH}_i$ y si está a menos de $\beta\cdot\mathrm{DBH}_i$.
  \item[Criterio de competencia.] Nueve índices de competencia, independientes y dependientes de la distancia, evaluados por el coeficiente de determinación ajustado de un modelo lineal del incremento de volumen en función del volumen y del índice.
  \item[Método.] Regresión lineal con optimización de $\alpha$ y $\beta$ mediante una búsqueda iterativa que maximiza el $R^{2}$ ajustado.
  \item[Resultados.] El modelo solo con volumen explicó 32.5\,\% de la variabilidad y el que incluye competencia, 64\,\%. El mejor índice fue el de Hegyi por altura y distancia ($R^{2}$ ajustado de 0.647), los índices que consideran la distancia superaron a los que no, y optimizar la distancia influyó más que optimizar el umbral de DBH. Los autores señalan que no pudieron evaluar la competencia interespecífica por la baja proporción de otras especies.
\end{description}

\subsubsection{Síntesis}

La \cref{tab:comparacion-competencia} resume las cinco referencias.

\begin{table*}[t]
  \centering
  \caption{Comparación de las referencias sobre reducción de competencia.}
  \label{tab:comparacion-competencia}
  \small
  \begin{tabularx}{\textwidth}{@{}>{\RaggedRight\arraybackslash}p{2.7cm}>{\RaggedRight\arraybackslash}p{3.1cm}*{2}{>{\RaggedRight\arraybackslash}X}@{}}
    \toprule
    Referencia & Dominio & Criterio de competencia o incompatibilidad & Método \\
    \midrule
    \parencite{elizalde2025} & Reforestación & Coeficiente por par de especies en posiciones vecinas & Programación entera y algoritmo genético \\
    \parencite{eisenblatter2002} & Telecomunicaciones & Valor de interferencia por par de transmisores (0 a 1) & Programación entera, heurísticas y cotas \\
    \parencite{akplogan2011} & Agricultura & Efecto del cultivo previo, compatibilidad con el suelo y agrupación espacial & WCSP con ramificación y acotamiento \\
    \parencite{liu2026} & Agricultura bajo incertidumbre & Matriz de interacción entre cultivos vecinos & Optimización robusta con escenarios Monte Carlo \\
    \parencite{kobal2025} & Manejo forestal & Índices de competencia según tamaño y distancia de vecinos & Regresión con optimización de umbrales \\
    \bottomrule
  \end{tabularx}
\end{table*}

Las referencias muestran tres formas de medir la competencia: coeficientes asociados a pares de elementos vecinos en un grafo \parencite{elizalde2025,eisenblatter2002,liu2026}, penalizaciones por el efecto del elemento previo o cercano \parencite{akplogan2011} e índices que combinan tamaño y distancia \parencite{kobal2025}. En cuanto a los métodos, los modelos exactos pierden escalabilidad al crecer el problema \parencite{elizalde2025,akplogan2011}, lo que justifica el uso de heurísticas y metaheurísticas \parencite{elizalde2025,eisenblatter2002}. Por otro lado, solo \textcite{liu2026} incorpora incertidumbre mediante escenarios Monte Carlo, aunque la aplica a rendimientos, precios y costos. Ninguna de las cinco estima la cantidad de plantas que se deben adquirir cuando se desconocen las plantas preexistentes, que es el segundo componente del reto y donde la simulación Monte Carlo resulta necesaria.

%==================================================
\subsection{Problemas de coloración en grafos}
\label{sec:coloracion}

\subsubsection{Coloración de vértices}

Sea $G=(V,E)$ un grafo no dirigido. Una coloración propia de vértices con $k$ colores es una función $c:V\to\{1,\dots,k\}$ tal que $c(u)\neq c(v)$ para toda arista $uv\in E$. El menor $k$ para el que existe una coloración propia es el número cromático $\chi(G)$. En el problema de asignación de frecuencias, cada vértice representa un transmisor, dos vértices se unen si sus señales interfieren cuando usan la misma frecuencia y los colores corresponden a frecuencias; el número mínimo de frecuencias necesarias es entonces el número cromático del grafo de conflicto \parencite{eisenblatter2002}. Los mismos autores recuerdan que la coloración de vértices es un problema NP-difícil.

Como la coloración propia no basta para modelar requisitos prácticos, existen generalizaciones: la T-coloración (distancias prohibidas entre colores de vértices vecinos), la coloración por listas (colores bloqueados por vértice), la coloración por conjuntos y los modelos que, en lugar de prohibir los conflictos, minimizan una penalización total de interferencia \parencite{eisenblatter2002}. Según \textcite{dewerra2015}, los modelos básicos de coloración de vértices y de aristas se han extendido para resolver problemas de programación (\textit{scheduling}) cada vez más complejos.

\subsubsection{Coloración de aristas}

Una coloración propia de aristas con $k$ colores es una función $f:E\to\{1,\dots,k\}$ tal que dos aristas que comparten un vértice reciben colores distintos. El menor $k$ posible es el índice cromático $\chi'(G)$. Según \textcite{dewerra2015}, las coloraciones de aristas se aplican a problemas de talleres abiertos (\textit{open shop}), horarios escolares y calendarios deportivos.

Una variante es la coloración fuerte de aristas adyacentes \parencite{zhang2002}: es una coloración propia de aristas en la que, además, para cada arista $uv$ los conjuntos de colores incidentes en $u$ y en $v$ son distintos, es decir $C(u)\neq C(v)$ con $C(u)=\{f(uv)\mid uv\in E\}$. El menor número de colores se denota $\chi'_{as}(G)$. Los autores muestran que $\chi'_{as}(G)\geq\Delta(G)$, donde $\Delta(G)$ es el grado máximo; que para un árbol vale $\Delta$ si no hay dos vértices de grado máximo adyacentes y $\Delta+1$ si los hay; y proponen la conjetura de que $\chi'_{as}(G)\leq\Delta(G)+2$ para todo grafo conexo con al menos seis vértices. También señalan que algunos problemas de redes se pueden convertir a esta clase de coloraciones.

\subsubsection{Representación del problema de reforestación mediante un grafo}

La distribución de especies en un polígono puede representarse con un grafo donde cada elemento del problema tiene un equivalente (\cref{tab:grafo-reforestacion}). Esta representación coincide con la de \textcite{elizalde2025}, que modelan el terreno como un grafo de posiciones de siembra y vecindades, y con la matriz de adyacencia de \textcite{liu2026}. La definición de vecindad puede variar (por ejemplo, solo vecinos horizontales y verticales o también diagonales) y determina la densidad del grafo.

\begin{table*}[t]
  \centering
  \caption{Representación propuesta del problema de reforestación como un grafo.}
  \label{tab:grafo-reforestacion}
  \small
  \begin{tabularx}{\textwidth}{@{}p{3.2cm}X@{}}
    \toprule
    Elemento del grafo & Equivalente en el problema de reforestación \\
    \midrule
    Vértices & Posiciones de siembra del polígono \\
    Aristas & Pares de posiciones vecinas entre las que puede haber competencia \\
    Colores & Especies a plantar \\
    Peso de arista & Nivel de competencia o incompatibilidad entre las especies asignadas a sus extremos \\
    \bottomrule
  \end{tabularx}
\end{table*}

\paragraph{Coloración de vértices como referencia} Una coloración propia exigiría que posiciones vecinas tengan especies distintas, pero no considera por sí sola la demanda por especie ni que la competencia entre pares de especies sea parcial. Por ello, el problema se parece más a las generalizaciones que minimizan una penalización total que a la coloración propia. Se observa además que el modelo de mínima interferencia de \textcite{eisenblatter2002} y el modelo de \textcite{elizalde2025} tienen la misma estructura: variables binarias de asignación (canal o especie por vértice), una restricción que asigna exactamente un valor a cada vértice y variables auxiliares que activan una penalización en cada arista cuando sus extremos toman valores incompatibles, con el objetivo de minimizar la suma. Esto convierte a la coloración de vértices con penalización en la referencia más cercana para desarrollar la estrategia de distribución de especies, y a las heurísticas de coloración (como DSATUR) en candidatas a construir soluciones iniciales.

\paragraph{Coloración de aristas} Las aristas del grafo (pares de posiciones vecinas) no tienen un equivalente natural de especie, por lo que esta coloración no sirve directamente para decidir qué especie va en cada posición. Una posibilidad, que aquí se plantea únicamente como hipótesis de trabajo, es emplearla para el orden de plantación, interpretando los colores como etapas o turnos en los que se atienden las interacciones entre posiciones vecinas, de forma análoga a los usos en programación de horarios señalados por \textcite{dewerra2015}.

\subsubsection{Conclusión de la sección}

Los problemas de coloración de vértices, y en particular sus generalizaciones con penalización por conflicto, ofrecen un marco adecuado y bien estudiado para la distribución de especies, y la literatura revisada indica que, al crecer el tamaño, conviene recurrir a heurísticas o metaheurísticas. La coloración de aristas aporta, en todo caso, una idea complementaria para el orden de plantación. La estimación de la cantidad de plantas a adquirir queda fuera de estos modelos y se aborda con las herramientas de probabilidad y simulación de la \cref{sec:incertidumbre}.

\section{Conceptos en la modelación con incertidumbre}

Consideramos los siguientes conceptos que serán necesarios para desarrollar un modelo que actúe en un proceso estocástico.


\subsection{Probabilidad basica}


\begin{definition}
    Una variable aleatoria $X$ es una función de un espacio muestral $\Omega=\{s_1,s_2,s_3,\dots\}$ hacia los números reales. Con una funcion de probabilidad $\mathbb{P}$ en $\Omega$, sea $\mathcal{X}=\{x_1,\dots,x_m\}$ el rango de $X$. Definimos la funcion de probabilidad $\mathbb{P}_X$ en $\mathcal{X}$ observando que $X=x_i$ si y solo si el resultado de un experimento aleatorio es $s_j\in\Omega\;\text{t.q}\;X(s_j)=x_i$. Entonces:
    
    \begin{align}
        \mathbb{P}_X(X=x_i) = \mathbb{P}(\{s_j\in\Omega\,:\,X(s_j)=x_i\})
    \end{align}\autocite{:berger2002}.
\end{definition}

Para toda variable aleatoria $X$ le asociamos una función de distribución acumulada.

\begin{definition}
    Una funcion de distribución acumulada $F_X(x)$ es definida como:

    \begin{align}
        F_X(x)=\mathbb{P}_X(X\leq x)\;\forall\;x
    \end{align}\autocite{:berger2002}.
\end{definition}


Una variable aleatoria puede estar asociada a una función de probabilidad acumulada \textit{cdf} que puede ser o una función de densidad de probabilidad \textit{pdf} o una función de masa de probabilidad \textit{pmf}, esto dependiedo si esta en un espacio continuo o discreto respectivamente.

\begin{definition}
    Una \textit{pmf} para una variable aleatoria discreta es

    \begin{align}
        f_X(x)=\mathbb{P}_X(X=x)\,\forall\,x.
    \end{align}
\end{definition}

\begin{definition}
    Una \textit{pdf} de una variable aleatoria continua satisface

    \begin{align}
        F_X(x)=\int_{-\infty}^{x}f_X(t)\mathbf{dt}\,\forall\,x
    \end{align}\autocite{:berger2002}.
\end{definition}


Entre las distribuciones de probabilidad que podrian.
\section{Identificación inicial del modelo de optimización}
\label{sec:modelo-optimizacion}

\subsection{Función objetivo}
\label{subsec:funcion-objetivo}

Hay más de un criterio razonable. Conviene identificarlos antes de comprometerse con uno solo.

\begin{itemize}
\item \textbf{Minimizar la competencia entre especies vecinas}, sumada sobre todos los pares de posiciones adyacentes de la retícula. Es el criterio que responde directamente a la parte de distribución espacial del reto.
\item \textbf{Minimizar el costo esperado de compra}, sumando el costo de cada especie por la cantidad que hay que adquirir. Responde a la parte de cantidad a comprar.
\item \textbf{Una combinación ponderada de ambas}, ya que la distribución espacial también puede influir en la supervivencia y, con ello, en la compra futura, y porque el reto permite mover las proporciones por especie dentro de un margen de $\pm$5\% a $\pm$10\%, lo que conecta ambas decisiones.
\end{itemize}

De estas, minimizar la competencia espacial es el objetivo más directamente pedido por el reto para la primera decisión, y se tomará como punto de partida; el costo se incorporará en cuanto se conozcan precios por especie y el alcance real del presupuesto de \$13{,}200.

\subsection{Restricciones}
\label{subsec:restricciones}

\Cref{tab:restricciones} enumera las restricciones identificadas hasta ahora y la condición del problema que representa cada una.

\begin{table*}[t]
\centering
\caption{Restricciones identificadas para el modelo de optimización.}
\label{tab:restricciones}
\small
\begin{tabular}{@{}p{0.22\textwidth}p{0.42\textwidth}p{0.28\textwidth}@{}}
\toprule
Restricción & Qué condición representa & Por qué se incluye \\
\midrule
Cantidad requerida por especie & El total de plantas de cada especie debe cumplir el plan (658 plantas/ha repartidas según las proporciones dadas, ajustadas dentro de un margen de $\pm$5\% a $\pm$10\%) & Es la meta explícita del reto; sin ella el modelo podría vaciar una especie completa \\
Capacidad o espacio disponible & El número de posiciones en la retícula de cada polígono es fijo, según su superficie y la densidad del plan & No se puede plantar más individuos de los que caben en el área disponible \\
Distancias mínimas entre plantas & La separación entre posiciones vecinas queda fija por el arreglo al tresbolillo ($\approx$4.19 m con el plan actual) & Es una condición del propio diseño de plantación que el reto pide respetar \\
Compatibilidad entre especies & Ciertos pares de especies no deberían quedar como vecinos inmediatos, o deberían evitarse en exceso & Es el criterio ecológico que sostiene la función objetivo de competencia \\
Disponibilidad de plantas & La cantidad comprada no puede superar lo que el vivero puede producir o entregar en la temporada de siembra & Vuelve factible en la práctica la cantidad que el modelo calcule, aunque el dato aún no se tiene \\
Requerimientos de agua & Las especies con mayor demanda de agua podrían necesitar restricciones adicionales de ubicación o densidad & Relevante en un matorral desértico, aunque no está cuantificado en los datos actuales \\
Existencia previa de plantas & Las posiciones donde ya hay una planta establecida no deberían recibir una nueva, y esas plantas cuentan hacia la meta del plan & Evita comprar y plantar de más donde ya se cumplió la meta \\
Presupuesto & El costo total de la compra no debe superar los \$13{,}200 disponibles & Límite económico explícito del reto, pendiente de precisar su alcance \\
\bottomrule
\end{tabular}
\end{table*}

Esta identificación es preliminar. No fija todavía coeficientes ni una formulación matemática completa (eso corresponde a los siguientes entregables), pero delimita qué debe entrar al modelo y por qué, y distingue las restricciones que ya se pueden cuantificar con los datos actuales (cantidad por especie, capacidad, distancias) de las que dependen de información que aún falta por obtener (compatibilidad, disponibilidad, agua, presupuesto).


\nocite{*}
\printbibliography

% --- Anexos ---
% Solo informacion complementaria citada desde el cuerpo: las tecnicas que se
% probaron y descartaron (\cref{sec:tecnicas-descartadas}) y los diccionarios
% de datos (\cref{sec:diccionario,sec:diccionario-transformado}). No cuentan
% contra el limite de paginas del cuerpo.
\appendix

\end{document}
